\section{引言}
\label{sec:introduction}

在UTXO账本中，付款花费的是更早交易创建的输出。若收款方必须等待该交易公开执行，付款链的每一跳都会增加一次确认延迟。设Alice向Bob支付100 CAL，Bob在Alice的交易执行之前就把该输出转给Carol。Bob需要可验证的证据，证明这笔价值有支撑。若Bob的付款先执行，账本便消费了一个尚不存在的输出，缺失的100 CAL需要一个明确的来源。无论Alice的交易先于Bob的到达、后于Bob的到达还是始终不到达，账目都必须正确。

最直接的替代方案是在每一跳都等待公开执行，第~\ref{sec:evaluation}节将其作为对照。基于证书的支付系统已把收款方可用性与公开处理分开：FastPay和Zef支持低延迟的证书化支付，Sui Lutris针对不同类型的状态访问结合广播与共识路径~\cite{fastpay,zef,lutris}。支付通道依靠通道关系及其流动性实现快速转账~\cite{blitz,donner}。Next要回答的是：当来源尚未执行时，原生账本如何为一笔已认证的付款提供资金，由注册组织而非预充值通道为未解决的来源担责。

四个难点相互耦合。证书在公开登记之前就已流通，因此从公开记录的缺口出发的核算会漏掉已存在的承诺。迟到的来源不能既向垫付价值的担保方付款，又向已花费该输出的收款方付款。客户端可能持有来源却永不提交。并且证书不是执行承诺：$\mathsf{IntentID}=\operatorname{Digest}(\texttt{INTENT},\mathit{Network},\mathit{Subject},\mathit{Nonce})$ 对同一网络上的同一付款人全局有效，因此恶意用户可以就两笔输入相互独立的付款从两个组织分别取得法定人数证书（QC），而公开执行对该Intent只接受一个TxID。证书因此界定了其发行方的负债，却不保证来源会执行或回款。

我们的洞见是：对缺失来源负责的一方，已由随输出流通的证书指明。Next把缺口记在该直接发行方名下，而不追溯祖先；用有备付支撑的签发额度约束每份已签证书的负债；并以一个无需改写已存历史的公开决定关闭义务。我们把设计组织为三项贡献。

\noindent\textbf{C1：面向一轮续花的分层担保架构。}成员在签名前原子地锁定输入并预留签发额度，一轮并行请求产生TXCer。收款方验证并原子记录该输出后即可请求后继付款；请求携带其直接输入的证书，因此对固定形状，其证据不随深度增长。组织内的复制与公开提交在后台进行，付款本金与担保方备付相互分离。

\noindent\textbf{C2：直接责任驱动的来源解耦。}若被认证的输入缺失，公开执行原子地消费它，并记录针对发行方的义务。后继所在区块携带来源的证书，原签发成员可据已存批准补交来源。期限届满后，公开决定扣减发行方备付，无需门限适配，迟到的来源再偿还该扣减。未执行的来源在签发额度中保留其全部CAL金额（含找零），部分批准同样占用额度。对于四成员组织和四验证者委员会（每个组织及委员会至多一个静态拜占庭成员），我们证明唯一消费，并证明每份证书的未结CAL负债无论是否已公开均有资金覆盖。对于有限、无环、来源闭合、无冲突且每笔付款都成功执行一次的付款集合，终态CAL持有量不依赖到达顺序与赔付时点；FUEL不在其内。我们不主张这些证明适用于一般委员会规模。

\noindent\textbf{C3：授权历史表示。}受限变色龙哈希命令把已赔付输入的资金引用替换为备付引用，同时保持所有者授权、后继引用和完整BlockID。永久决定是经济依据，成功执行的Revision与Task授权确切的表示；仅哈希匹配不构成授权。经济闭合既不等待表示，也不依赖表示。其接口收益是：已验证并执行公共前缀$H$的诚实副本，可在一次读事务中按原坐标$(h,j,k)$读出该输入由哪份备付垫付。代价是每个被修复块额外的存储记录和一次本地维护。

在一台M4 Max主机上，以冻结版本做的机制内对照给出以下结果。在成员同步提交、钱包存储不同步的配置下，100跳链到最后一次证书到账的中位时间为4.557\,s，而每跳等待公开执行的对照为64.476\,s；快速链的全部297个后继中，首次发送都早于其父付款最早的委员会提交。十四服务混合负载使用组织代付FUEL，每种处理三轮，共完成84,192笔付款，客户端存在背压，见第~\ref{sec:evaluation}节；原签发成员解决了补交处理中全部24个被扣留来源，适配服务全部暂停时赔付与回款仍然完成。在21,600次不含对象缓存的本地热页查询中，相对优化决定索引，表示使点读P50降低7.44--12.32\%，主要来自块解码；代价是每块新增0.28--2.19\,MiB逻辑记录，以及一次43--214\,ms的本地维护，不含网络与共识。